% ===================================================================
\chapter{Organización del proyecto}

\section{Mapa de carpetas}
\begin{archivo}[label=estructura]
pulso/
├── README.md               qué es el proyecto, resultados y cómo usarlo
├── CHANGELOG.md            historial de versiones
├── pyproject.toml          ficha del paquete: nombre, versión, dependencias, comando pulso
├── uv.lock                 versiones exactas de todas las dependencias
├── .python-version         versión de Python (3.12)
├── Dockerfile              receta de la imagen de la aplicación
├── .dockerignore           lo que no entra en la imagen
├── .gitignore              lo que git no guarda (entorno, cachés, temporales)
├── .github/workflows/ci.yml   integración continua: ruff, pytest y Docker
├── .streamlit/config.toml  tema y opciones del servidor de la aplicación
├── configs/                TODOS los parámetros, en YAML
├── src/pulso/              el código del modelo (paquete de Python)
├── app/                    la aplicación web (Streamlit)
├── data/                   datos generados por los experimentos
├── models/                 modelos entrenados y versionados
├── figures/                figuras de las tareas 1 a 6
├── notebooks/              análisis y figuras de cada tarea
├── tests/                  pruebas automáticas
└── docs/                   documentación, fuentes LaTeX y PDF
\end{archivo}

La regla del proyecto es separar responsabilidades: los \textbf{parámetros} viven en
\ruta{configs/}, el \textbf{cálculo} en \ruta{src/pulso/}, la \textbf{presentación} en
\ruta{app/} y los \textbf{resultados} en \ruta{data/}, \ruta{models/} y
\ruta{figures/}. Ningún número del modelo está escrito en el código.

\section{El código: \texttt{src/pulso/}}

\subsection*{Núcleo del modelo (tareas 1 a 4)}
\begin{longtable}{L{0.24\textwidth}L{0.68\textwidth}}
\toprule
\textbf{Archivo} & \textbf{Qué hace} \\
\midrule
\endhead
\texttt{networks.py} & Genera redes de pagos: Erdős–Rényi, Barabási–Albert y
núcleo-periferia, con montos lognormales escalados a media 1. \\
\texttt{model.py} & Define los parámetros del modelo (\texttt{Params}), los colchones
de impago y el lado derecho de las ecuaciones de déficit y fase, incluido el
acoplamiento multicapa. \\
\texttt{integrate.py} & Integra las ecuaciones en el tiempo (Runge–Kutta 4, ruido de
Euler–Maruyama), aplica choques y detecta impagos y cascadas. \\
\texttt{spectral.py} & Análisis espectral: radio espectral, índice $R$, estabilidad
$\sigma$, tiempo de recuperación, $\lambda_2$, comunidades y ranking de entidades. \\
\texttt{cascades.py} & Modelos clásicos de contagio para comparar: Eisenberg–Noe,
DebtRank y Watts. \\
\texttt{metrics.py} & Indicadores de cada simulación: fracción en impago, liquidez no
pagada, primer contagio, colapso e importancia sistémica. \\
\bottomrule
\end{longtable}

\subsection*{Experimentos y predicción (tarea 5)}
\begin{longtable}{L{0.24\textwidth}L{0.68\textwidth}}
\toprule
\textbf{Archivo} & \textbf{Qué hace} \\
\midrule
\endhead
\texttt{experiments.py} & Barrido de escenarios: muestrea redes y parámetros, simula y
guarda las tablas con su esquema validado. \\
\texttt{ml/train.py} & Entrena y valida (GroupKFold por red) la regresión logística,
XGBoost y la GCN; guarda métricas. \\
\texttt{ml/gnn.py} & Red neuronal sobre grafos (GCN) con conexión residual. \\
\texttt{ml/pinn.py} & Red neuronal informada por la física para estimar parámetros. \\
\texttt{ml/predict.py} & Predicción instantánea para una red nueva, explicación SHAP y
verificación por simulación. \\
\texttt{ml/registry.py} & Guarda y carga modelos con versión y metadatos. \\
\bottomrule
\end{longtable}

\subsection*{Cartera de un banco (tarea 6)}
\begin{longtable}{L{0.24\textwidth}L{0.68\textwidth}}
\toprule
\textbf{Archivo} & \textbf{Qué hace} \\
\midrule
\endhead
\texttt{world/schema.py} & Constantes (tipos de cliente, capas) y esquemas de las tablas
de la cartera. \\
\texttt{world/entities.py} & Genera personas y empresas con sus variables y su
probabilidad de impago verdadera. \\
\texttt{world/layers.py} & Genera las relaciones de nómina, proveedores y
transferencias. \\
\texttt{world/core.py} & Reúne todo en un objeto \texttt{World}; lo guarda y lo carga. \\
\texttt{world/transactions.py} & Convierte la simulación en un registro de
transacciones como el de un banco. \\
\texttt{credit/engine.py} & Parámetros y simulación del modelo multicapa, colchones según
la PD e impagos idiosincráticos. \\
\texttt{credit/estimate.py} & Estima la red, las fases y la coherencia a partir de las
transacciones. \\
\texttt{credit/features.py} & Calcula las variables de cada cliente en tres niveles:
clásico, red estática y red dinámica. \\
\texttt{credit/pipeline.py} & Une generación, simulación, transacciones y variables para
una cartera. \\
\texttt{credit/evaluate.py} & Validación cruzada por cartera, métricas, diferencias
pareadas e intervalos de confianza. \\
\texttt{credit/experiment.py} & Experimento completo de la tarea 6: carteras, regímenes,
resultados y modelos finales. \\
\texttt{credit/service.py} & Funciones que usa la aplicación en las páginas 8 y 9. \\
\bottomrule
\end{longtable}

\subsection*{Utilidades}
\begin{longtable}{L{0.24\textwidth}L{0.68\textwidth}}
\toprule
\textbf{Archivo} & \textbf{Qué hace} \\
\midrule
\endhead
\texttt{config.py} & Lee los YAML de \ruta{configs/} y resuelve las carpetas de datos y
modelos (admite variables de entorno). \\
\texttt{schema.py} & Valida columnas, tipos y rangos de las tablas, y prohíbe usar las
columnas de verdad \texttt{truth\_*}. \\
\texttt{service.py} & Funciones que usa la aplicación en las páginas 1 a 7. \\
\texttt{viz.py} & Gráficas interactivas (Plotly) de la aplicación. \\
\texttt{cli.py} & El comando \texttt{pulso} (capítulo~\ref{cap:cli}). \\
\bottomrule
\end{longtable}

\section{Las demás carpetas}
\begin{longtable}{L{0.3\textwidth}L{0.62\textwidth}}
\toprule
\textbf{Carpeta o archivo} & \textbf{Contenido} \\
\midrule
\endhead
\ruta{app/Inicio.py} & Página de inicio de la aplicación. \\
\ruta{app/common.py} & Estado compartido y cachés de la aplicación. \\
\ruta{app/pages/1\_…9\_*.py} & Una página por archivo (capítulo~\ref{cap:app}). \\
\ruta{configs/*.yaml} & Parámetros (capítulo~\ref{cap:config}). \\
\ruta{data/*.parquet}, \texttt{networks.npz} & Conjunto de datos de la tarea 5; ficha en
\ruta{data/README.md}. \\
\ruta{data/credit/} & Resultados del experimento de la tarea 6. \\
\ruta{models/} & Modelos de la tarea 5 por nombre y versión; \ruta{models/credit/v1/}
para la tarea 6. \\
\ruta{figures/} & Figuras de los notebooks y de la aplicación. \\
\ruta{notebooks/01…06} & Un notebook por tarea, del 01 (red y déficit) al 06 (PINN). \\
\ruta{tests/test\_fase*.py} & Pruebas por fase; \texttt{test\_app.py} prueba que cada
página de la aplicación se ejecute sin errores. \\
\ruta{docs/Presentaciones/} & PDF finales: v1, v2, v3, manuales y plan de estudio. \\
\ruta{docs/latex/} & Fuentes LaTeX, generador de figuras (\texttt{figuras.py}) y
\texttt{compilar.sh}. \\
\ruta{docs/diseno/} & Especificación de diseño de la tarea 6. \\
\ruta{docs/bibliografia.md} & Referencias por tarea. \\
\bottomrule
\end{longtable}

\begin{aviso}
No se guardan en git: el entorno (\ruta{.venv/}), las cachés de Python y la tabla
grande de clientes \ruta{data/credit/clientes.parquet}, que se regenera con
\texttt{uv run pulso credit configs/credit.yaml}.
\end{aviso}
